\documentclass[12pt,a4paper]{article}
\usepackage[utf8]{inputenc}
\usepackage[T1]{fontenc}
\usepackage{enumitem}
\usepackage[margin=2.5cm]{geometry}
\usepackage{parskip}
\usepackage{amsmath}
\usepackage{amssymb}
\usepackage{xcolor}

\title{midterm --- Final Exam --- Answer Key}
\date{}

\begin{document}

\maketitle

{\large \textbf{\textcolor{red}{ANSWER KEY --- Not for Distribution}}}

\vspace{0.5em}



\noindent
\textbf{Total points:} 4
\hfill
\textbf{Number of questions:} 4

\vspace{1em}
\hrule
\vspace{1.5em}


\noindent
\textbf{Question 1 (1.00 pts).}~
What is a key difference between a classical bit and a qubit?


\begin{enumerate}[label=\textbf{\Alph*.}]

  \item A classical bit is more like a coin spinning in the air

  \item A qubit is always in exactly one of two positions: 0 or 1

  \item \textbf{\textcolor{teal}{[CORRECT]}}  A qubit is more like a coin spinning in the air --- until you catch it, the outcome is undecided

  \item A classical bit is used for quantum computing

\end{enumerate}



\textit{\small Explanation: The key difference between a classical bit and a qubit is that a classical bit is always in one of two positions (0 or 1), while a qubit is like a coin spinning in the air, with an undecided outcome until measured.}



\vspace{0.8em}
\noindent\rule{\linewidth}{0.2pt}
\vspace{0.3em}


\noindent
\textbf{Question 2 (1.00 pts).}~
The Hadamard gate can create a superposition from a basis state.


\textbf{Answer:} \textcolor{teal}{True}



\textit{\small Explanation: The Hadamard gate is specifically designed to create superpositions from basis states, allowing for the exploration of multiple possibilities simultaneously.}



\vspace{0.8em}
\noindent\rule{\linewidth}{0.2pt}
\vspace{0.3em}


\noindent
\textbf{Question 3 (1.00 pts).}~
What is the key difference between the correlation of entangled qubits and the correlation of matching socks?


\textbf{Model Answer:} \textcolor{teal}{The correlation of entangled qubits is created at the moment of measurement, whereas the correlation of matching socks is decided when they are packed.}



\textit{\small Explanation: Entangled qubits have no definite value until measured; Matching socks have a predetermined correlation}



\vspace{0.8em}
\noindent\rule{\linewidth}{0.2pt}
\vspace{0.3em}


\noindent
\textbf{Question 4 (1.00 pts).}~
Explain the relationship between decoherence and the scalability of quantum computing, including the role of T1 and T2 in this context.


\textbf{Key Points:} \textcolor{teal}{Introduction to decoherence, explanation of T1 and T2, discussion of scalability challenges, and analysis of how different physical platforms address these challenges.}



\textit{\small Explanation: Consider the physical platforms used by companies like IBM Quantum, Google Quantum AI, and Xanadu, and how decoherence affects their ability to scale up the number of qubits.}





\end{document}
